\section*{Capítulo \ref{ch:b3:extension-strike-slip} --- Tectônica extensional e transcorrente}

\begin{solution}{exo:b3:extension-strike-slip:1}
As camadas mergulham $50 - 35 = \ang{15}$, afastando-se das falhas; $\beta =
\sin\ang{50}/\sin\ang{35} = 0.766/0.574 = 1.34$.
\end{solution}

\begin{solution}{exo:b3:extension-strike-slip:2}
Escalonamento dextral para a direita: bacia. Escalonamento dextral para a esquerda: elevação. Escalonamento
sinistral para a esquerda: bacia. Escalonamento sinistral para a direita: elevação.
\end{solution}

\begin{solution}{exo:b3:extension-strike-slip:3}
$105/20 = \qty{5.25}{km/Myr} = \qty{5.25}{mm/yr}$, próximo de
\qty{6}{mm/yr}: a taxa não precisou ser constante ao longo de \qty{20}{Myr}.
\end{solution}

\begin{solution}{exo:b3:extension-strike-slip:4}
\qty{200}{\percent} significa $\Delta L = 2L_0$: $L_0 = \qty{117.5}{km}$ e $L
= \qty{352.5}{km}$. No norte, $\Delta L = 0.5L_0$: $L_0 = \qty{470}{km}$
e $L = \qty{705}{km}$.
\end{solution}

\begin{solution}{exo:b3:extension-strike-slip:5}
$563/20 = 28$ a $563/15 = \qty{38}{mm/yr}$. O deslocamento é um valor mínimo, e
a deriva medida por levantamento topográfico através da zona de falhas também inclui o movimento em
falhas vizinhas do sistema e a deformação elástica acumulada entre
terremotos.
\end{solution}

\begin{solution}{exo:b3:extension-strike-slip:6}
$t = 90 - \text{mergulho}$. Em \ang{60}: $(1 + 0.6\times1.732)/(1 - 0.6\times
0.577) = 2.039/0.654 = 3.12$; em \ang{50}: $1.715/0.497 = 3.45$; em
\ang{40}: $1.504/0.285 = 5.28$. Travamento em $\arctan0.6 = \ang{31.0}$.
\end{solution}

\begin{solution}{exo:b3:extension-strike-slip:7}
$\arctan0.4 = \ang{21.8}$ e $\arctan0.2 = \ang{11.3}$. Uma falha com mergulho de
\ang{15} só desliza se $\mu < \tan\ang{15} = 0.27$.
\end{solution}

\begin{solution}{exo:b3:extension-strike-slip:8}
$2.5\times4.0 = \qty{10.0}{km^2}$ e $2.5\times(4.0 - 1.5) =
\qty{6.25}{km^2}$.
\end{solution}

\begin{solution}{exo:b3:extension-strike-slip:9}
$132/1.19 = \qty{110.9}{km}$, estirada em \qty{21.1}{km}, o que leva
$21.1/6 = \qty{3.5}{Myr}$.
\end{solution}

\begin{solution}{exo:b3:extension-strike-slip:10}
$3.0/\sin\ang{60} = \qty{3.464}{km}$ antes e $3.0/\sin\ang{45} =
\qty{4.243}{km}$ depois: \qty{0.779}{km} por bloco ($\beta = 1.22$). As camadas
mergulham $60 - 45 = \ang{15}$. Cada bloco gira como um todo: a rotação que
reduz o mergulho das falhas que mergulham, digamos, para leste, faz uma linha horizontal
no bloco subir para leste, e as camadas mergulham para oeste, em sentido oposto às
falhas.
\end{solution}

\begin{solution}{exo:b3:extension-strike-slip:11}
Em uma zona sinistral leste--oeste, $\sigma_1$ tem rumo \ang{045}. Cisalhamentos R a
$\phi/2 = \ang{20}$ da zona, no lado de $\sigma_1$: direção \ang{070},
sinistral. R$'$ a $90 - 20 = \ang{70}$: direção \ang{020}, dextral. Fraturas de tração
paralelas a $\sigma_1$: direção \ang{045}.
\end{solution}

\begin{solution}{exo:b3:extension-strike-slip:12}
Para $\mu = 0.6$, as falhas se formam a \ang{60.5} e travam a \ang{31.0}: $\beta =
\sin\ang{60.5}/\sin\ang{31.0} = 0.870/0.515 = 1.69$. Para $\mu = 0.85$: de
\ang{65.2} a \ang{40.4}, $\beta = 0.908/0.648 = 1.40$. Falhas com atrito menor giram
mais antes de travar: rochas fracas se estendem mais antes que novas
falhas precisem se formar.
\end{solution}

\begin{solution}{pb:b3:extension-strike-slip:1}
\textbf{1.} As camadas eram horizontais e giraram com os blocos; as
falhas giraram do mesmo ângulo, logo $\theta_0 = 42 + 18 = \ang{60}$,
o mergulho de Anderson para $\phi = \ang{30}$.

\textbf{2.} $\beta = \sin\ang{60}/\sin\ang{42} = 0.8660/0.6691 = 1.294$.

\textbf{3.} $24.0/1.294 = \qty{18.5}{km}$; o domínio foi estendido em
\qty{5.5}{km} (\qty{5.46}{km}).

\textbf{4.} $t = \ang{48}$: $(1 + 0.6\times0.900)/(1 - 0.6\times1.111) =
1.540/0.334 = 4.62$, contra $3.12$ para falhas novas a \ang{60.5}: quase
metade a mais.

\textbf{5.} Em $\arctan0.6 = \ang{31.0}$; $\beta = \sin\ang{60}/\sin\ang{31}
= 1.68$ no máximo. Antes disso, falhas novas a cerca de \ang{60} cortariam
os blocos, e uma segunda geração de dominós assumiria o controle.

\textbf{6.} Cada coluna do bloco superior desce pela queda da falha sob ela sobre o
afastamento horizontal; a falha é mais íngreme perto da superfície e se achata
à medida que se afasta dela, de modo que as colunas próximas à falha descem mais e as mais distantes
menos: as camadas se inclinam em direção à falha.

\textbf{7.} $D = A/H = 16.0/3.2 = \qty{5.0}{km}$.

\textbf{8.} $D$ é proporcional a $A$: $\qty{5.0}{km} \pm \qty{0.5}{km}$.

\textbf{9.} $3.2\times(5.0 - 2.0) = \qty{9.6}{km^2}$.

\textbf{10.} O afastamento horizontal, \qty{3.2}{km}.

\textbf{11.} O preenchimento engrossa em direção à falha, em forma de cunha, e suas
camadas mais antigas estão mais inclinadas do que as mais novas: cada camada foi depositada enquanto
a charneira de rollover e o espaço junto à falha continuavam crescendo.

\textbf{12.} $z = P/(\rho g) = \num{5.5e8}/(2800\times9.8) = \qty{20.0}{km}$.

\textbf{13.} $20.0/(24 - 16) = \qty{2.5}{km/Myr} = \qty{2.5}{mm/yr}$.

\textbf{14.} $20.0/\sin\ang{30} = \qty{40}{km}$ de deslizamento, e $40\cos\ang{30} =
\qty{34.6}{km}$ de extensão horizontal.

\textbf{15.} Não: \ang{15} está abaixo do mergulho de bloqueio de \ang{31}. Ou a
falha deslizou enquanto era mais íngreme e foi depois rotacionada (\omterm{def:b3:extension-strike-slip:core-complex}{charneira rolante}), ou
ela era fraca, com $\mu < 0.27$ devido a argilas ou fluidos sob alta pressão.

\textbf{16.} Por $45 - 15 = \ang{30}$: à medida que a placa superior se afastava, o
bloco inferior descarregado subiu por isostasia e se curvou, rotacionando a falha através da
charneira.

\textbf{17.} \qty{34.6}{km} frente aos \qty{8.7}{km} do semigráben (questão
23): quatro vezes mais. \omterm{def:b3:extension-strike-slip:core-complex}{Complexos de núcleo metamórfico} ocorrem onde a extensão é maior, como no
centro da Basin and Range.

\textbf{18.} Uma falha sinistral com degrau para a esquerda: \omterm{def:b3:extension-strike-slip:bends}{curva de alívio}, abre-se uma bacia.

\textbf{19.} $\qty{5}{mm/yr}\times\qty{1.2}{Myr} = \qty{6.0}{km}$.

\textbf{20.} $4.0\times6.0 = \qty{24}{km^2}$; $3.0 + 6.0 = \qty{9.0}{km}$ de comprimento;
$9.0/4.0 = 2.25$.

\textbf{21.} $132/10 = 13$ a $132/7 = 19$: a bacia do Mar Morto é muito mais
alongada, com deslocamento muitas vezes maior que sua largura.

\textbf{22.} Uma \omterm{def:b3:extension-strike-slip:bends}{curva de restrição}: uma crista empurrada para cima entre os ramos; em
seção, uma \omterm{def:b3:extension-strike-slip:flower}{estrutura em flor} positiva, com a falha se abrindo para cima em ramos
com componente reversa, semelhante ao que tem de 100 a \qty{300}{m} de largura na Dead
Sea Transform.

\textbf{23.} $3.2 + 5.46 = \qty{8.66}{km}$: o semigráben de Merrow foi
estendido em \textbf{\qty[text-series-to-math=true]{8.7}{km}}.
\end{solution}
